\subsection{What one extent lookup costs}
\label{sec:fiemap}

Table~\ref{tab:fiemap} times the fork's own lookup function (\texttt{stat},
\texttt{open}, FIEMAP, \texttt{close}) over files with a controlled number of extents
on four file systems, warm and cold, and counts the device reads per file. The extent
counts in the table are measured, not intended: our punch-and-rewrite procedure gave
exactly one extent per block on all four file systems.

\begin{table}[t]
\caption{Extent lookup per file (median of five passes over 100--1{,}000 files), and
device reads per file in a cold pass. One thread unless stated.}
\label{tab:fiemap}
\small
\inputtable{../results/tab_fiemap.tex}
\end{table}

\paragraph{Warm, it is microseconds.} With the cache warm, a lookup on ext4 or XFS
costs 16--25\,\textmu s for up to eight extents (30--44\,\textmu s on Btrfs) and grows
with the extent count, because the ioctl copies every extent out (about 0.6\,ms for
512 extents). This is the normal case for FragPicker: the analysis runs right after
the traced application has read these files, so their inodes and extent blocks are in
memory. F2FS is the exception: about 0.4\,ms per lookup even for one extent; we did
not find where that time goes.

\paragraph{Cold, the cost is the extent tree, not the inode.} With the cache cold,
files whose extents fit in the inode (four on ext4, more on XFS) still need almost no
device reads: 0.01 reads per file on ext4 and 0.04 on XFS, because inode-table
readahead brings in the inodes of hundreds of neighbouring files at once. Once the
extents spill into a tree block, every file needs one read of its own (1.0 per file
on both, two to three for 512 extents), and the lookup cost rises to 0.4--1.6\,ms on
our loop device. F2FS reads one block per file at every extent count, the file's
inode, which F2FS stores in a block of its own.

\paragraph{Threads helped only where the wait is outside FIEMAP.} On ext4 and XFS,
four and sixteen threads were \emph{slower} per file than one in every row
(0.21--0.96$\times$), even in the rows where each file needs its own device read. On
F2FS they were 1.4--1.6$\times$ faster in every row, and on Btrfs with 128 or more
extents 1.1--1.2$\times$. The difference matches where the wait happens: F2FS's one
read per file comes in through the \texttt{stat} and \texttt{open} calls, which release
the interpreter lock, while the read that remains on ext4 and XFS, the extent-tree
block, happens inside FIEMAP, where the fork's request keeps the lock held
(\S\ref{sec:gil}). We did not time the four calls separately, so this attribution is
an inference from the read counts.
